\chapter{Deep Learning Model Implementation}
This chapter comprehensively describes the architecture of the Transformer-CNN-LSTM deep learning model, which consumes the unified 3D point cloud generated by the dual-radar fusion pipeline to perform binary fall classification.

\section{Unified Point Cloud as Model Input}
The output of the fusion pipeline is a rich 3D cloud with enhanced spatial density, completely overcoming the single-sensor occlusion problem. The DL model ingests this unified point cloud, utilizing the $z$-axis (height), radial velocity, and SNR as core features. Rapid downward velocity combined with a collapsing $z$-spread serves as the fundamental physical signature of a fall.

\section{Hybrid Model Architecture Rationale}
Analyzing time-series radar data requires understanding complex spatio-temporal dynamics. A fall event spans multiple seconds and can easily be confused with similar vertical motions, such as sitting down quickly. A hybrid architecture was designed to fuse three distinct neural modeling paradigms:
\begin{itemize}
\item \textbf{Transformer Encoder:} Employs multi-head self-attention to capture global temporal dependencies, establishing relationships between distant frames (e.g., the pre-fall stance versus the post-fall resting state) simultaneously.
\item \textbf{1D Convolutional Neural Network (CNN):} Sweeps across the temporal axis to extract local, high-frequency kinetic patterns (e.g., sudden jerks or sudden accelerations) while reducing the sequence length through max-pooling.
\item \textbf{Long Short-Term Memory (LSTM):} Consumes the CNN's output to strictly model the sequential, time-ordered evolution of the features.
\end{itemize}

\section{Enhanced Architecture Layout}
The model ingests a batch of sliding windows, shaped \texttt{(Batch, 40, 20)}, where the 20 features are extracted from the unified point cloud. 

\subsection{Projection and Attention Blocks}
First, a Linear Projection layer expands the input into 64 dimensions. A sinusoidal Positional Encoding is added to retain sequential awareness. The data then flows through two sequential \texttt{TransformerEncoderLayer} blocks. Each block features Multi-Head Self-Attention (4 heads, $d_{model}=64$) and a Feed-Forward Network with an inner dimension of 128. GELU activations and Layer Normalization are utilized for training stability.

\subsection{Convolutional Feature Compression}
The temporal output is transposed for CNN consumption. The CNN block features two stacked Conv1D layers (64 $\rightarrow$ 128 $\rightarrow$ 128) with kernel sizes of 3. Interspersed with BatchNorm1d and Dropout (0.4), the block concludes with a MaxPool1d layer of kernel 2, effectively halving the time dimension from 40 down to 20 steps, filtering high-frequency noise.

\subsection{Recurrent Memory and Classification}
The compressed sequence is fed into a Bidirectional LSTM (hidden size 128, providing 256 total output features). Rather than naively taking the final hidden state, the network employs a Learnable Attention Pooling mechanism over the LSTM outputs. Finally, a Multi-Layer Perceptron (MLP) classification head equipped with heavy Dropout layers maps the pooled features to the 2 output logits representing the "Fall" and "No-Fall" probabilities.

\section{Training Configuration}
The binary model was optimized using the AdamW optimizer with a learning rate of 0.001 and weight decay of 0.001. To combat the inherent class imbalance, a Class-Weighted CrossEntropyLoss function was employed alongside a label smoothing factor of 0.1 to prevent overconfidence and improve generalization on sparse radar datasets. Gradient clipping (max norm = 1.0) was strictly enforced to stabilize the LSTM gradients.
